\section{模型敏感性与适用边界}
数值验证检验的是离散算法是否可靠求解已建立的方程。界面关系、物性与省略机制引入的不确定性属于另一层面，需要用明确条件下的比较讨论。本节情景仅用于判断哪些假设值得关注，不替换第一问主结果。

\subsection{环境插值与表面交换}
在同一1280区间径向网格上，保持其他条件不变，分别采用分段线性与保形分段三次插值（PCHIP）处理附件1，所得温度在$1800\times21$个输出采样点上的最大差约0.0042\unit{K}。该差异大于已检验的数值离散误差，表明输出保留四位小数并不意味着对输入处理也具有同等稳定性。正式结果采用分段线性输入，以保留测量节点且避免区间过冲。

在同一1280区间网格的确定性情景计算中，将$h_m$分别降低和提高20\%，1800\unit{s}的表面含水率分别为1.6636、1.3759\unit{kg/kg}，同网格基准为1.5102\unit{kg/kg}。在所检情景中，增大传质系数使末时表面含水率进一步降低。这一扰动幅度是人为设置的敏感性区间，不是根据实验标定获得的参数置信区间；它提示真实交换系数及其浓度基准，比继续追求更多数值小数位更值得检验。

\subsection{潜热与界面平衡的一致性}
若额外将$820\unit{kg/m^3}$解释为初始湿体密度，则固定骨架干密度为$\rho_d=\frac{820}{1+2.55}=230.9859\unit{kg/m^3}$。再假定主质量模型的全部等效失水均在侧面蒸发，并取水汽化潜热$\lambda=2.45\unit{MJ/kg}$作量级估计\cite{fao}，由平均含水率变化推得前1800\unit{s}的条件性失水约0.01861\unit{kg}，相应潜热约45.59\unit{kJ}，为主模型显热增量约4.801\unit{kJ}的9.50倍。因此，不能依据“预热阶段”名称或未给潜热参数，声称相变吸热一定很小。

然而，保留上述失水通量$j_w=\rho_d h_m(C_s-C_\infty)$并直接在表面能量方程增加$\lambda j_w$，也不能自动形成自洽的相变模型。该单向耦合诊断计算得到1800\unit{s}表面温度约11.90\celsius。若进一步将环境量解释为空气湿度比$W$，并假设总压$p=101.325\unit{kPa}$，由湿空气关系\cite{ashrae}及饱和蒸气压关系\cite{fao}，
\begin{equation}
p_v=\frac{pW}{0.621945+W},\qquad
p_{\mathrm{sat}}(T)=0.6108\exp\!\left(\frac{17.27T}{T+237.3}\right)\unit{kPa},
\label{eq:psychrometric}
\end{equation}
其中式\eqref{eq:psychrometric}中的$T$以摄氏度计，得末时环境水蒸气分压约5.1156\unit{kPa}、露点约33.3008\celsius；11.90\celsius{}处纯水饱和蒸气压仅约1.3933\unit{kPa}。在这些附加条件及通常的水活度不大于1的解释下，原定持续向外净蒸发不再与气相驱动力相符。

该诊断说明，完整相变描述必须同时建立空气状态、表面水活度或吸附关系、质量通量与能量收支。题目未给这些材料关系，因此本文保留显热与等效扩散的基础模型，并将诊断低温视为闭合不一致的信号，而非真实药材温度预测；这一信号也不反向证明潜热可以忽略。

\subsection{几何收缩与辐射情景}
附件2给出半径在前1800\unit{s}由2.000\unit{cm}减小至1.873\unit{cm}，相对减少6.35\%；若长度固定，对应体积减少约12.30\%。这两项是几何变化量，不能直接解释成温度或含水率预测误差。第一问固定圆柱模型用于分问比较，第四问需对移动边界和守恒关系作一致处理。

在额外假设壁温等于烘房气温、周围为大黑体腔体且视角因子为1时，加入辐射的表面条件为
\begin{equation}
-kT_r(R,t)=h_T(T_s-T_\infty)
+\varepsilon\sigma\big[(T_s+273.15)^4-(T_\infty+273.15)^4\big],
\label{eq:radiation}
\end{equation}
其中$\sigma$为Stefan--Boltzmann常数，四次方项采用绝对温度。取$\varepsilon=1$的情景与同网格无辐射情景比较，1800\unit{s}轴心和表面温度分别增加0.6044\unit{K}、0.8095\unit{K}。更换时间积分法后的最大差小于$2.70\times10^{-8}\unit{K}$，辐射情景320至640径向区间的对照差小于$6.12\times10^{-6}\unit{K}$，足以分辨该情景效应。

由于真实壁温、发射率和接触条件未知，上述增量只能说明模型对额外辐射机制存在可分辨的敏感性，不能称为实际辐射误差或一般上界。主模型尚无内部温湿实测响应用于校准，故本稿的可靠性结论限定为所列假设、题给参数及输入方式下的数值解。
